Sigui la funció $f(x)=\dfrac{ax^2+x+b}{x^2+1}$.

\begin{apartats}

\apartat{1,25}
Calculeu els valors dels paràmetres $a$ i $b$ si sabem que la gràfica de la funció $f$ té un extrem relatiu en $x=-1$
i passa pel punt $P=\left(-2,\frac{13}{5}\right)$.

\begin{solucio}
La funció derivada de $f(x)$ és $f'(x)=\dfrac{-x^2+(2a-2b)x+1}{\left(x^2+1\right)^2}$. Imposem la condició d'extrem
relatiu: $f'(-1)=0\rightarrow\frac{-1+2b-2a+1}{4}=0\rightarrow a=b$. Imposem que la funció passi pel punt
$P=\left(-2,\frac{13}{5}\right)$:
\[
f(-2)=\frac{13}{5}\;\rightarrow\;\frac{4a-2+b}{5}=\frac{13}{5}\;\rightarrow\;4a+b=15 .
\]
D'on es dedueix que $a=b=3$.

\textit{Pauta oficial:} 0,25 pel càlcul de la funció derivada, 0,25 per imposar la condició d'extrem i trobar una
primera condició per als paràmetres, 0,25 per imposar que passi pel punt donat i trobar una segona condició, i 0,5
per trobar els valors dels paràmetres $a$ i $b$.
\end{solucio}

\apartat{1,25}
Per al cas $a=b$, calculeu i classifiqueu els extrems relatius de la funció.

\begin{solucio}
Per a $a=b$, la derivada és $f'(x)=\dfrac{1-x^2}{\left(x^2+1\right)^2}$. Resolem $f'(x)=0$ per trobar els possibles
extrems relatius, i obtenim $x=1$ i $x=-1$. Com que $f'(-2)<0$, $f'(0)>0$ i $f'(2)<0$, deduïm que la funció té un
mínim relatiu en $x=-1$ i un màxim relatiu en $x=1$.

\textit{Pauta oficial:} 0,25 per donar la funció derivada en el cas demanat, 0,5 per trobar els dos candidats a
extrems i 0,5 per donar i justificar on fa un mínim i on fa un màxim.
\end{solucio}

\end{apartats}
